% ============================================================
% CHAPTER 7: DISCUSSION
% Citations reuse verified entries in bibliography.bib only.
% ============================================================

This section brings together the empirical results from \autoref{sec:results}
and draws out their interpretation, design implications, and scope.
Section~\ref{sec:results:synthesis} reads all four motifs as faces of a single
amplifier mechanism.
Section~\ref{sec:discussion:interpretation} develops the interpretation of that
mechanism, its two operating regimes, and the status of null results.
Section~\ref{sec:discussion:design} translates the findings into empirically-grounded
design principles and the groundwork for a diagnostic suite.

\subsection{The Amplifier Mechanism: Reading the Motifs Together}
\label{sec:results:synthesis}

The preceding sections measured four motifs and a high-repetition distribution.
Read together, they do not describe four separate phenomena but four views of a
single underlying mechanism.
This subsection states that mechanism explicitly, shows how each measurement is a
signature of it, and separates what the data establish from what they do not.

\paragraph{The convergent finding.}
Across every analysis in the preceding section, the property that most strongly governs
the group's final answer is the composition of the votes \emph{before any debate
takes place}.
The high-repetition experiment shows this observationally: whether the
ground-truth option is present at $t = 0$ moves per-task accuracy from near zero
(0.4--6.9\,\%) to 36--58\,\%, and the seeded-initialisation experiment shows it
causally: placing a single correct agent into an otherwise-wrong opening profile
raises pooled accuracy roughly thirteen-fold (1.8\,\% to 23.7\,\%), in all 34
tasks tested (Sections~\ref{sec:results:evid} and~\ref{sec:results:seeded}).
The self-reinforcement analysis supplies the mechanism by which the opening
profile is converted into the outcome: SR is universal
($p_{\mathrm{SR}} \in [0.73, 0.88]$ in every cell), it is the convergence engine
(task-level $\rho(\bar\beta, \text{rounds}) \in [-0.89, -0.73]$ everywhere), and
its accuracy effect is a pure amplification, positive when the correct option
is already present ($r_{pb} = +0.20$ GPQA / $+0.15$ HB) and negative when it is
absent ($-0.11$ / $-0.07$), cancelling to near zero in the margin.
In short, the debate reliably \emph{amplifies} whichever position the opening
distribution favours; it only rarely \emph{generates} a position no agent
initially held (de-novo recovery in 15 of 34 seeded tasks, but at just 1.8\,\%
of repetitions).
We state this as the central empirical claim of the thesis: within the system we
tested, deliberation functions predominantly as a selective amplifier of its
initial vote distribution rather than as a mechanism that reasons a new answer
into existence.
This claim is not novel in kind. \citet{oh2026entrenchment} name homogenised
dynamics that amplify the majority view regardless of correctness as one of two
root causes of belief entrenchment, and \citet{estornell2024multillm} derive the
same conclusion theoretically from response homogeneity. What our data add is
measurement: a per-motif decomposition of how the amplification is realised,
conditioned on correctness throughout, and a causal test of the input on which
it operates.

\paragraph{How each motif is a signature of the same mechanism.}
The remaining three motifs describe \emph{how} the amplification plays out, not a
separate process.
\emph{Bistability} is the cross-repetition face of amplification: because the outcome is
selected by the (stochastic) opening profile, repeated repetitions of the same task land
in different basins, so multi-attractor coexistence is common
($N_{\mathrm{eff}} \geq 1.5$ in 78.6\,\% of cells) and is a stable property of the
task rather than of the configuration (cross-config $\rho = 0.78$--$0.91$).
The within-task accuracy relationship is dataset-specific: on GPQA higher
$N_{\mathrm{eff}}$ predicts lower accuracy ($\rho = -0.57$), because competing
attractors create wrong-answer basins in a setting where agents usually start
correct; on HiddenBench it predicts higher accuracy ($\rho = +0.41$), because
bistability gives the minority-correct faction a competing stable basin in a
setting where the dominant attractor is wrong.
The negative coupling between bistability and self-reinforcement
(task-level $\rho(N_{\mathrm{eff}}, \bar\beta) = -0.79$ on GPQA) is exactly what the
amplifier picture predicts: strong within-basin amplification means fast lock-in
to one attractor per repetition, whereas high attractor multiplicity dilutes the SR
signal across competing basins.
\emph{Dominance} is the within-repetition face: amplification is typically driven by a
small, concentrated cascade ($D$ correlates $\approx 0.87$ on GPQA and $\approx 0.76$ on HiddenBench at the task level with the top source's share and \emph{negatively} with total conversion volume), i.e.\ one agent tips a
near-balanced start rather than the group converging by diffuse attrition.
Configurations with higher dominance also require more rounds to converge
($\rho = +0.55$ GPQA, $+0.73$ HiddenBench, within-task), because subordinate agents
resist before capitulating.
\emph{Limit cycles} are the boundary case in which amplification fails to
terminate: when the opening profile is symmetric between two defensible options
and the star hub sits in the minority, each flip reconstructs the opposing
majority and the group never locks in, running to the hard cap.
Finally, the effective topology analysis supplies a within-run complement: even with
a fixed FC structure, the emergent influence star, where one agent captures most
conversions, achieves near-ceiling accuracy because the correct agent tends to win
the influence competition.
This is the amplifier operating at a finer grain: the debate amplifies whichever
agent the opening competition selects, and diverse initialisation raises the
probability that competition is won by the correct one.

\paragraph{The one axis that matters, and the one that does not.}
Two negative results sharpen the picture.
First, the \emph{structural} design axes we varied, communication topology and
memory window, do not change \emph{whether} the system reaches the correct
answer: topology has no accuracy effect at the task level (the confidence
interval for the fc--star gap contains zero on both benchmarks), the memory
window has none that survives multiple-comparison correction (a single
HiddenBench/fc cell is nominally significant, $q = 0.064$), and bistability is
invariant to both.
Their influence is confined to the \emph{dynamics}: star topology is slower,
is the sole source of genuine limit cycles, and mechanically concentrates
influence; larger memory windows speed HiddenBench convergence slightly.
The structural topology parameter is distinct from the \emph{emergent} influence
topology: conditional on FC runs, whether an effective star forms does predict
accuracy substantially, but this emergence is driven by the opening vote
composition rather than by the structural configuration, so the same conclusion
holds at every level of analysis, the opening vote is the dominant lever.
Second, the system's own internal signals are not indicators of correctness:
confidence rises by essentially the same amount on correct and incorrect repetitions
($+1.63$ vs.\ $+1.64$ on GPQA, with incorrect rising very slightly more), and on HiddenBench early unanimity is
anti-correlated with accuracy (round-2 stops reach only 2.5\,\%).
What \emph{does} move accuracy, the opening vote composition, is set before
any of these knobs act.

\paragraph{Scope of the claim.}
These statements are bounded by what we varied and what we held fixed.
We manipulated topology, memory window, task difficulty, benchmark, and (in the
seeded experiment) the opening vote composition; we did \emph{not} vary the base
model, the number of agents ($N = 4$), agent heterogeneity, or the cooperative
framing.
The amplifier characterisation is therefore established \emph{for a homogeneous,
cooperative, 4-agent debate on one model family across two benchmarks}, with
one axis (initial composition) tested causally and the rest observationally.
Whether the same mechanism dominates for heterogeneous models, larger groups, or
adversarial protocols is an open question we return to in
\autoref{sec:limitations}.
The value of the motif lens is that each of these
signatures is separately measurable from ordinary interaction logs, so a
practitioner can check which regime a given system is in rather than having to
assume it.


\subsection{Interpretation of Findings}
\label{sec:discussion:interpretation}

The preceding section argued that all four motifs are signatures of the same
amplifier mechanism.
This section unpacks what that implies: how SR should be read as a dynamical
process rather than a benefit, how the two benchmarks place the system in
qualitatively different operating regimes, why the structural null results are
informative, and what the robustness check does and does not show.

\paragraph{Self-reinforcement as an amplifier, not a benefit.}
The single most robust behavioural finding is that self-reinforcement is universal
and is the engine of convergence, yet is blind to correctness: it is present in
every configuration and predicts fewer rounds everywhere, but its relationship to
accuracy reverses sign depending on whether the correct option is present in the
opening votes. This is the empirical content of describing the system as an
amplifier: reinforcement strengthens whatever plurality exists at the outset,
which helps when that plurality is right and hurts when it is wrong. The reading
aligns with, and gives a dynamical mechanism for, the critical literature that
finds debate often fails to beat simple baselines \citep{liu2025stopovervaluing}
and that majority voting alone accounts for much of the apparent gain
\citep{zhang2025debateorvote}: if the process mostly amplifies the initial vote,
then its expected value over a population of repetitions is close to that of the opening
vote itself. It is also consistent with the anti-Bayesian confidence drift
documented by \citet{prasad2025certain}, which we observe directly as
outcome-independent confidence gain.
The reversal of the predicted $W$-slope ordering on GPQA (§\ref{sec:results:sr})
further supports a peer-conformity reading: shorter memory reduces each agent's
exposure to peer uncertainty, allowing confidence to ratchet more freely regardless
of correctness.

One comparison we deliberately do not make is against chain-of-thought or
self-consistency: our question is not whether debate is the best inference-time
method but what \emph{happens inside} it, so the fair reference point is the round-0
majority vote of the same four agents, and our ``gain'' figures are stated relative
to that pre-deliberation baseline rather than to an external single-agent method.

\paragraph{Two regimes: refinement and correction.}
The two benchmarks are not merely easy and hard versions of one task; they place
the system in two qualitatively different regimes, and the motif signatures track
the difference. On GPQA the agents usually begin with the correct option present
(oracle $70.6\,\%$), so amplification is mostly \emph{refinement}: SR consolidates
a frequently-correct start, the group ends above the $25\,\%$ chance baseline
($42.2\,\%$), and higher influence concentration is neutral-to-mildly-helpful. On
HiddenBench the agents begin systematically wrong, majority voting at $t=0$
scores below random ($7.6\,\%$) because errors are correlated, the LLM analogue of
biased information sampling in human hidden-profile groups
\citep{stasser1985pooling, li2026hiddenbench}; yet the same amplifier mechanism
operates here on a biased starting distribution.
The $+18.5$\,pp gain over the majority baseline reflects those repetitions where a
correct initial signal is amplified, not a systematic reasoning correction: the
oracle bound ($46.3\,\%$) confirms that the correct answer is available in roughly
half of repetitions, and amplification of those openings drives the recovery.
The motif differences follow: bistability is deeper on HiddenBench and, crucially,
its relationship to accuracy reverses sign between datasets, helping on HiddenBench
(where the dominant attractor is often wrong and a competing basin is beneficial)
but hurting on GPQA (where competing attractors create wrong-answer basins that
the system can fall into); the group's final accuracy is no better than chance,
limit cycles concentrate there, and influence concentration cuts both ways, helping
in the fully-connected topology but \emph{lowering} accuracy when forced through a
single hub (star). The same mechanism (amplification of the opening profile) thus
produces opposite-looking phenomenology depending on the reliability of that
profile, which is precisely why a practitioner cannot read off system health from
accuracy alone.

A complementary asymmetry appears at the individual level.
On GPQA, agents who vote correctly at $t=0$ are marginally more confident than
those who do not (mean $7.39$ vs.\ $7.25$; point-biserial $r = +0.047$, $p < 0.001$):
correctness and confidence align, so initially-correct agents carry both a numerical
majority and a confidence advantage into deliberation.
On HiddenBench the direction reverses ($r = -0.092$, $p < 0.001$; $7.20$ correct
vs.\ $7.39$ wrong): an agent whose single private fact happens to mislead them is
more confident than one whose fact points to the correct answer.
The initially-correct minority thus faces a double obstacle---structural outvoting by
three peers and a confidence deficit---which accounts for the smaller
informed-agent influence advantage on HiddenBench ($+1$\,pp) versus GPQA ($+15$\,pp).

\paragraph{Structural axes shape dynamics, not correctness.}
Topology and memory window, the two axes most often tuned in the debate
literature \citep{li2024sparsetopology, kaesberg2025voting}, had no accuracy
effect that survives multiple-comparison correction, but clear effects on the
dynamics: the star
topology is slower, concentrates influence roughly threefold, and is the only
configuration that produces genuine (hub-driven) limit cycles, while larger memory
windows modestly speed HiddenBench convergence. Both halves of this dissociation
are corroborated externally. On the null side, network type has no significant
main effect on coordination in an eight-topology factorial
\citep{mehdizadeh2026topologymemory} and communication density does not lift
accuracy above a non-communicative baseline at any level tested
\citep{bertalanic2026costofconsensus}. On the dynamics side, the same
eight-topology study finds a large memory-by-topology interaction on settling
time, confirming that these axes move the trajectory even where they do not
move the outcome. This dissociation between
dynamical and epistemic effects is itself a finding: it means these knobs are
levers on \emph{how} the group deliberates (cost, worst-case runtime,
entrenchment) but not on \emph{whether} it is right, and it explains why
topology-tuning results in the literature are mixed \citep{yang2025revisiting}.
The key qualification is that the structural topology parameter is distinct from the
\emph{emergent} influence topology.
Within FC runs, effective star formation, where one agent captures most of the
conversion credit and the correct agent tends to win that competition, predicts
above-average accuracy; this structure is inaccessible by assigning the star
parameter but arises naturally from diverse opening positions.
Prescribing a structural hub fixes the center before the debate can select it, and
on hard questions the randomly assigned center started correct only 15.0\,\% of
the time versus 85.5\,\% for a debate-elected hub.
The mechanism is direct: within effective-star runs, when the hub votes correctly at
$t=0$ group accuracy reaches $86.1\,\%$ on GPQA and $92.4\,\%$ on HiddenBench;
when it votes incorrectly, accuracy collapses to $15.6\,\%$ and $10.7\,\%$
respectively.
The hub's opening vote---a function of its randomly assigned private
information---is the single largest predictor of group outcome in the star
configuration, which explains why prescribing the star parameter is unreliable
without also controlling which agent occupies the center.

\paragraph{Null and near-null results as findings.}
Several results are informative precisely because they are negative, and the study
was designed to treat them as such (\autoref{sec:intro:rqs}). Genuine limit cycles
are essentially absent in the fully-connected topology; bistability is invariant to
topology and memory; memory window does not move accuracy; and the marginal
SR--accuracy correlation is near zero. The last is not a true null but a
cancellation of two opposite conditional effects, which we could only see by
conditioning on ground-truth presence, a reminder that an aggregate null can
hide structured, offsetting effects. Reporting these against principled null models
(shuffle surrogates, permutation tests, sign tests) is what separates ``no effect''
from ``not enough power to see one,'' and we have tried to keep that distinction
explicit throughout.

\paragraph{Robustness to the system prompt.}
The agent persona used throughout includes an explicit anti-conformity
instruction, so we can ask whether the motif structure is an artefact of that
particular wording. A paired comparison against an earlier repetition with a milder
persona, on the same HiddenBench questions and configurations
(Appendix~\ref{app:sysprompt}), shows that the more critical prompt changes the
\emph{dynamics}, faster fully-connected convergence, fewer cap-hits, and a higher
unanimous-convergence rate, but does not change accuracy (0 of 6 configurations)
and does not suppress self-reinforcement (0 of 6). The apparent shift toward more
multistable labels is a by-product of more repetitions converging rather than a genuine
change in attractor multiplicity. The amplifier mechanism thus survives a
deliberate attempt to make agents less conformist, which is mild evidence that it
is a property of the interaction structure rather than of a single prompt. This is
a single HiddenBench comparison, and is not evidence of robustness across models or
group sizes, which we did not test.
\citet{bertalanic2026costofconsensus} argue on structural grounds that
prompt-level remedies are unlikely to resolve these dynamics: adversarial
prompting risks widening the oracle gap by increasing the rate at which valid
reasoning is attacked, and distinct prompts across a homogeneous parameter space
produce stylistic variance without decoupling the underlying error distributions.
Our null result on accuracy is consistent with that expectation.


\subsection{From Dynamics to Diagnostics: Design Implications}
\label{sec:discussion:design}

The motivation for this thesis was that outcome metrics tell a practitioner
\emph{whether} a configuration worked but never \emph{why}, and give no way to
tell a healthy deliberation from a pathological one \citep{miehling2025agentic}.
The \suitename{} detectors are a first, partial answer: each is computed from
ordinary interaction logs and each reports on a distinct aspect of the
deliberative dynamics. We do not claim to deliver a validated diagnostic tool.
What the evidence supports is more modest and, we think, more honest: a set of
\emph{empirically-grounded design principles}, tiered by how strongly our data
support them, together with the observation that the very quantities used to
detect each motif can serve as the instruments a future diagnostic suite would
need. This subsection states the principles and that mapping.

\paragraph{A tiered reading of the evidence.}
We separate the principles by evidential strength so that the confident claims are
not diluted by the speculative ones.
\emph{Tier 1 (strong: universal across our cells, and for the first, causal).}
(1)~The opening vote composition is the dominant lever on the outcome; investment
in the process that generates round-0 votes dominates tuning of the debate itself.
(2)~The system amplifies far more readily than it generates, so it is best used as
a consensus filter over candidate answers rather than as a mechanism expected to
discover an answer no agent holds. (3)~The system's own confidence and unanimity
are not signals of correctness and must not be used to gate trust
(corroborated independently by \citealt{bertalanic2026costofconsensus}, who
report near-zero confidence--correctness correlations and show that peer
exchange further degrades confidence discrimination).
\emph{Tier 2 (directional: consistent in sign, but dataset-dependent in size, so
``verify on your own setup'').} (4)~Communication topology is a lever on
efficiency and worst-case behaviour, not accuracy; the fully-connected topology is
the safer default and the star should be used only when a designated aggregator is
wanted. (5)~A hard round cap with explicit ``unresolved'' handling is necessary,
because only some topologies admit non-terminating debates and every observed limit
cycle ran to the cap
(failure to recognise termination conditions accounts for 12.4\% of annotated
multi-agent failures in deployed frameworks; \citealt{cemri2025mast}). (6)~Decisive leadership helps when all agents can see one
another (fully-connected) but \emph{hurts} when influence is forced through a
single hub (on HiddenBench/star, higher dominance lowers accuracy), so its value
depends both on the topology and on whether the leader is more likely correct.
(7)~Diverse initialisation on multi-option tasks raises the probability of
effective star formation from under $1\,\%$ to $23\,\%$; since the effective hub
disproportionately starts correct, this is a low-cost lever for above-average
accuracy on questions with $\geq\!4$ answer options
(consistent with enforced perspective diversity as an entrenchment remedy;
\citealt{oh2026entrenchment}).
(8)~Persistent non-updating agents are a vulnerability: a single devil's advocate
reduces accuracy from $63.7\,\%$ to $10.0\,\%$ even when outnumbered two-to-one,
so non-updating behaviour should be monitored and handled explicitly
(the convention-formation analogue is minority-driven norm change;
\citealt{ashery2025conventions}; see also \citealt{yao2025sycophancy} and
\citealt{pitre2025consensagent} for mitigation).
\emph{Tier 3 (diagnostic / minor knobs).} (9)~Bistability is a stable property of
the task, so tasks can be profiled offline, and stochastic tasks, where
convergence is effectively random and accuracy tends to be lowest, can be flagged
for routing elsewhere.
(10)~The memory window is a minor efficiency knob, not an accuracy lever.
(11)~Prompt changes (e.g.\ anti-conformity instructions) buy dynamical robustness,
not accuracy.

\paragraph{Each detector is an instrument.}
The practical bridge from these principles to a future tool is that every
principle above is checkable with a quantity we already compute
(\autoref{tab:design:instruments}). This is what we mean by ``groundwork for a
diagnostic suite'' rather than the suite itself: the measurements exist and have
been validated against null models on 30{,}000 repetitions; what remains, packaging
them into an online monitor, and validating the recommendations on models and group
sizes beyond ours, is future work, not a claim we make here.

\begin{table}[H]
\centering

\small
\begin{tabular}{p{4.3cm}p{3.5cm}p{5.2cm}}
\toprule
Principle & Observable (this thesis) & How a builder checks it \\
\midrule
Opening composition is the master lever (T1)
  & $P(\text{GT present})$, initial plurality; seeded contrast
  & Log the round-0 vote profile; relate to outcome across repetitions \\
\addlinespace
Amplifies, does not generate (T1)
  & De-novo recovery rate
  & Fraction of repetitions whose final answer was absent at $t=0$ \\
\addlinespace
Confidence / unanimity $\neq$ quality (T1)
  & Outcome-independent confidence gain; early-stop accuracy
  & Check self-confidence calibration against held-out truth \\
\addlinespace
Topology sets dynamics, not accuracy (T2)
  & Rounds, cap-rate, dominance $D$ by topology
  & Compare efficiency / worst-case across topologies, not accuracy \\
\addlinespace
Cap + unresolved-flagging is required (T2)
  & RQA determinism + period guard
  & Flag any repetition reaching the round cap as unresolved \\
\addlinespace
Profile tasks; route stochastic ones (T3)
  & $N_{\mathrm{eff}}$, Cram\'er's $V$, cell label
  & Pre-classify tasks offline; treat stochastic tasks with caution \\
\addlinespace
Diverse init raises P(effective star) (T2)
  & Initial vote entropy; effective topology type
  & On $\geq\!4$-option tasks, seed distinct starting positions to raise
    P(star) toward $23\,\%$ \\
\addlinespace
Non-updating agents are a vulnerability (T2)
  & Per-agent vote entropy across rounds
  & Flag agents with near-zero vote entropy; treat as suspect \\
\bottomrule
\end{tabular}
\caption[Design principles mapped to diagnostic measurements]{Mapping from design principle to the measurement that checks it. Every
         observable is computed in this thesis from ordinary interaction logs;
         the right column states how a practitioner would apply it to their own
         system. The suite itself is not built here, this is the set of
         instruments it would use.}
\label{tab:design:instruments}
\end{table}

\paragraph{Relation to enterprise deployment.}
The three Tier-1 principles are, in effect, safety principles for deploying such
systems in consequential settings. A system that confidently reaches unanimity on
a wrong answer, that cannot manufacture a correct answer absent from its inputs,
and whose outcome is largely fixed before deliberation begins, is one whose
reliability must be established \emph{upstream} (in the quality and diversity of
the agents that produce the opening positions) and \emph{externally} (through
verification that does not trust the system's own confidence), not through longer
or more elaborate debate. This is a deflationary message, but a useful one: it
directs effort toward the levers that our data show actually move the outcome and
away from those that do not. We emphasise once more that these recommendations are
grounded in a single model family, four agents, and two benchmarks, and should be
treated as hypotheses to be re-checked, using exactly the instruments in
\autoref{tab:design:instruments}, on any new deployment.
